\section{Tiered Execution}
\label{sec:tiered}

RaS continuity and tiered execution are related but separable claims. The continuity hypothesis can be tested using the same high-capability model for reasoning and editing. Tiered execution introduces a second question: after high-capability reasoning has produced a bounded work package, can routine stateful execution be delegated to a lower-cost, local, or more constrained worker without materially reducing engineering quality?

The high-capability reasoner is best reserved, under the proposed architecture, for operations such as architecture, decomposition, difficult causal reasoning, complex debugging, review, difficult trade-off analysis, interpretation of evidence, and work-package creation. An execution worker may instead perform repository edits, compilation, test execution, formatting, linting, Git operations, local-environment interaction, and tightly constrained environment actions where permitted.

\input{tables/reasoner-executor}

Let $N_R$ be the number of high-capability reasoning operations, $N_E$ the number of execution operations, $c_R$ the average cost of a high-capability operation, $c_E$ the average cost of an execution operation, and $C_O$ orchestration overhead. A simplified monolithic model is:
\begin{equation}
C_M=(N_R+N_E)c_R.
\end{equation}
Tiered execution is:
\begin{equation}
C_T=N_Rc_R+N_Ec_E+C_O.
\end{equation}
It is cheaper under this theoretical model when:
\begin{equation}
N_E(c_R-c_E)>C_O.
\label{eq:tiered}
\end{equation}

Equation~\ref{eq:tiered} is theoretical. It does not assume that $c_E<c_R$ in all deployments, that execution quality is equivalent, or that orchestration is cheap. A nominally inexpensive executor can become more expensive in outcome-normalised terms if it repeatedly misunderstands the work package, introduces regressions, exhausts retries, or forces the high-capability reasoner to re-enter.

Tiered execution can therefore fail through at least four mechanisms. First, the work package may omit a crucial semantic constraint. Second, the executor may lack enough capability to handle an unexpected repository state. Third, the executor may have insufficient tool access to reproduce the required evidence. Fourth, repeated hand-offs may create latency and context overhead larger than any saving.

Condition D in the broader research programme evaluates this architecture. P0 may initially focus primarily on persistent control, forced reset, and semantic-state ablation because those conditions test the central continuity proposition more directly. Tiered execution should not be allowed to confound the first continuity result.

The architecture nevertheless matters because it highlights a systems principle: the component that performs the most sophisticated reasoning does not necessarily need to perform the most machine operations. If continuity is external, reasoning capability, execution state, and operational privilege can be independently scoped.
